% General Relativity: Einstein Field Equations and Geodesic Analysis
% Topics: Schwarzschild metric, orbital precession, gravitational waves, cosmology
% Style: Theoretical physics report with computational analysis

\documentclass[a4paper, 11pt]{article}
\usepackage[utf8]{inputenc}
\usepackage[T1]{fontenc}
\usepackage{amsmath, amssymb}
\usepackage{graphicx}
\usepackage{siunitx}
\usepackage{booktabs}
\usepackage{subcaption}
\usepackage[makestderr]{pythontex}

% Theorem environments
\newtheorem{definition}{Definition}[section]
\newtheorem{theorem}{Theorem}[section]
\newtheorem{example}{Example}[section]
\newtheorem{remark}{Remark}[section]

\title{General Relativity: Computational Analysis of Curved Spacetime and Gravitational Phenomena}
\author{Theoretical Physics Laboratory}
\date{\today}

\begin{document}
\maketitle

\begin{abstract}
This report presents a comprehensive computational analysis of Einstein's general relativity,
exploring the curvature of spacetime and its observational consequences. We examine the
Schwarzschild metric for spherically symmetric spacetimes, compute geodesic trajectories
to analyze orbital precession (including Mercury's anomalous perihelion advance of 43
arcseconds per century), model gravitational wave signals from binary mergers, and
investigate cosmological solutions via the Friedmann-Lema\^{i}tre-Robertson-Walker metric.
Computational methods demonstrate how mass-energy curves spacetime geometry, producing
testable predictions from black hole event horizons to the large-scale structure of the universe.
\end{abstract}

\section{Introduction}

General relativity revolutionized our understanding of gravity by describing it not as a force,
but as the curvature of four-dimensional spacetime. The theory's central insight is that
mass-energy determines the geometry of spacetime, which in turn governs the motion of
particles and light through geodesic equations.

\begin{definition}[Einstein Field Equations]
The relationship between spacetime curvature and mass-energy is expressed as:
\begin{equation}
G_{\mu\nu} + \Lambda g_{\mu\nu} = \frac{8\pi G}{c^4} T_{\mu\nu}
\end{equation}
where $G_{\mu\nu} = R_{\mu\nu} - \frac{1}{2}Rg_{\mu\nu}$ is the Einstein tensor, $R_{\mu\nu}$
is the Ricci curvature tensor, $R$ is the Ricci scalar, $g_{\mu\nu}$ is the metric tensor,
$\Lambda$ is the cosmological constant, and $T_{\mu\nu}$ is the stress-energy tensor.
\end{definition}

\section{Theoretical Framework}

\subsection{Schwarzschild Metric}

\begin{theorem}[Schwarzschild Solution]
For a spherically symmetric, non-rotating mass $M$ in vacuum, the metric takes the form:
\begin{equation}
ds^2 = -\left(1 - \frac{r_s}{r}\right)c^2dt^2 + \left(1 - \frac{r_s}{r}\right)^{-1}dr^2 + r^2(d\theta^2 + \sin^2\theta\,d\phi^2)
\end{equation}
where $r_s = 2GM/c^2$ is the Schwarzschild radius (event horizon for black holes).
\end{theorem}

The metric coefficients reveal profound physics. The temporal component $g_{tt} = -(1-r_s/r)$
vanishes at $r=r_s$, indicating gravitational time dilation becomes infinite at the event
horizon. The radial component $g_{rr} = (1-r_s/r)^{-1}$ diverges at the same location,
demonstrating coordinate singularity (though spacetime itself remains regular in proper coordinates).

\subsection{Geodesic Equations}

\begin{definition}[Geodesic Motion]
Particles in free-fall follow geodesics, curves that extremize proper time:
\begin{equation}
\frac{d^2x^\mu}{d\tau^2} + \Gamma^\mu_{\alpha\beta}\frac{dx^\alpha}{d\tau}\frac{dx^\beta}{d\tau} = 0
\end{equation}
where $\tau$ is proper time and $\Gamma^\mu_{\alpha\beta}$ are Christoffel symbols encoding curvature.
\end{definition}

For the Schwarzschild metric, the effective potential for radial motion is:
\begin{equation}
V_{eff}(r) = -\frac{GM}{r} + \frac{L^2}{2r^2} - \frac{GML^2}{c^2r^3}
\end{equation}

The final term represents a general relativistic correction absent in Newtonian gravity,
producing orbital precession. For Mercury, this yields the famous 43 arcseconds per century
perihelion advance that confirmed Einstein's theory in 1915.

\subsection{Gravitational Waves}

\begin{theorem}[Linearized Gravity]
For weak gravitational fields, $g_{\mu\nu} = \eta_{\mu\nu} + h_{\mu\nu}$ where $|h_{\mu\nu}| \ll 1$,
the wave equation becomes:
\begin{equation}
\Box \bar{h}_{\mu\nu} = -\frac{16\pi G}{c^4}T_{\mu\nu}
\end{equation}
where $\bar{h}_{\mu\nu} = h_{\mu\nu} - \frac{1}{2}\eta_{\mu\nu}h$ and $\Box$ is the d'Alembertian.
\end{theorem}

Gravitational waves propagate at the speed of light, carrying two polarization states ($+$ and $\times$).
The strain amplitude from a binary merger scales as $h \sim (GM/c^2r)(v/c)^2$, where $M$ is the
chirp mass and $r$ is the distance to the source. LIGO's 2015 detection of GW150914 measured
strain amplitudes $h \sim 10^{-21}$, confirming this prediction.

\subsection{Cosmology: FLRW Metric}

\begin{definition}[Friedmann-Lema\^{i}tre-Robertson-Walker Metric]
For a homogeneous, isotropic universe:
\begin{equation}
ds^2 = -c^2dt^2 + a^2(t)\left[\frac{dr^2}{1-kr^2} + r^2(d\theta^2 + \sin^2\theta\,d\phi^2)\right]
\end{equation}
where $a(t)$ is the scale factor and $k \in \{-1, 0, +1\}$ determines spatial curvature.
\end{definition}

The Einstein equations reduce to the Friedmann equations:
\begin{align}
H^2 &= \left(\frac{\dot{a}}{a}\right)^2 = \frac{8\pi G}{3}\rho - \frac{kc^2}{a^2} + \frac{\Lambda c^2}{3} \\
\frac{\ddot{a}}{a} &= -\frac{4\pi G}{3}\left(\rho + \frac{3p}{c^2}\right) + \frac{\Lambda c^2}{3}
\end{align}

These equations govern cosmic expansion, with the cosmological constant $\Lambda$ driving
accelerated expansion observed via Type Ia supernovae (1998 Nobel Prize discovery).

\section{Computational Analysis}

\begin{pycode}
import numpy as np
import matplotlib.pyplot as plt
from scipy.integrate import odeint, solve_ivp
from scipy.optimize import fsolve

np.random.seed(42)

# Physical constants
G = 6.67430e-11  # m^3 kg^-1 s^-2
c = 299792458    # m/s
M_sun = 1.989e30 # kg
au = 1.496e11    # meters (astronomical unit)
year = 365.25 * 24 * 3600  # seconds

# Schwarzschild radius function
def schwarzschild_radius(mass):
    """Calculate Schwarzschild radius in meters"""
    return 2 * G * mass / c**2

# Schwarzschild metric coefficients
def metric_coefficients(r, M):
    """Return g_tt and g_rr for Schwarzschild metric"""
    r_s = schwarzschild_radius(M)
    if r <= r_s:
        return 0, np.inf
    g_tt = -(1 - r_s/r)
    g_rr = 1 / (1 - r_s/r)
    return g_tt, g_rr

# Effective potential for orbital motion
def effective_potential(r, M, L, relativistic=True):
    """Effective potential for radial motion in Schwarzschild metric"""
    r_s = schwarzschild_radius(M)
    V_newton = -G*M/r + L**2/(2*r**2)
    if relativistic:
        V_gr = -G*M*L**2/(c**2 * r**3)
        return V_newton + V_gr
    return V_newton

# Geodesic equations for orbital motion
def geodesic_equations(t, y, M):
    """
    Geodesic equations in Schwarzschild spacetime
    y = [r, phi, p_r, p_phi]
    """
    r, phi, p_r, p_phi = y
    r_s = schwarzschild_radius(M)

    if r <= r_s * 1.01:  # Near horizon cutoff
        return [0, 0, 0, 0]

    # Energy and angular momentum (constants of motion)
    E = 1.0  # Normalized energy
    L = p_phi  # Angular momentum

    # Geodesic equations
    dr_dt = p_r * (1 - r_s/r)
    dphi_dt = L / r**2

    # Radial acceleration
    dp_r_dt = L**2 / r**3 - G*M / r**2 + 3*G*M*L**2/(c**2 * r**4)
    dp_phi_dt = 0  # Angular momentum conserved

    return [dr_dt, dphi_dt, dp_r_dt, dp_phi_dt]

# Mercury orbital parameters
M_mercury = 3.285e23  # kg
a_mercury = 57.9e9    # semi-major axis (m)
e_mercury = 0.2056    # eccentricity
period_mercury = 87.97 * 24 * 3600  # orbital period (s)

# Calculate perihelion precession
def perihelion_precession_rate(a, e, M_central):
    """
    Calculate perihelion precession in arcseconds per orbit
    a: semi-major axis
    e: eccentricity
    M_central: central mass
    """
    # GR prediction
    delta_phi = 6 * np.pi * G * M_central / (c**2 * a * (1 - e**2))
    # Convert to arcseconds per orbit
    arcsec_per_orbit = delta_phi * 180 / np.pi * 3600
    return arcsec_per_orbit

# Gravitational wave from binary merger
def gravitational_wave_strain(t, M_chirp, distance, f_initial):
    """
    Simplified chirp waveform for inspiraling binary
    M_chirp: chirp mass in solar masses
    distance: luminosity distance in Mpc
    f_initial: initial frequency in Hz
    """
    M_chirp_kg = M_chirp * M_sun
    distance_m = distance * 3.086e22  # Convert Mpc to meters

    # Frequency evolution (leading order)
    tau = t - t[-1]  # Time to coalescence
    tau[tau >= 0] = -1e-6  # Avoid singularity

    f_t = f_initial * (1 - 256/5 * (np.pi * f_initial)**(8/3) *
                       (G * M_chirp_kg / c**3)**(5/3) * tau)**(-3/8)
    f_t[f_t < 0] = 0

    # Strain amplitude
    h_0 = (4 / distance_m) * (G * M_chirp_kg / c**2)**(5/3) * (np.pi * f_t / c)**(2/3)
    h_plus = h_0 * np.cos(2 * np.pi * f_t * t)

    return h_plus, f_t

# FLRW cosmology - scale factor evolution
def friedmann_equations(a, t, Omega_m, Omega_lambda):
    """
    Friedmann equation for flat universe (k=0)
    """
    H_0 = 2.2e-18  # Hubble constant in s^-1 (roughly 67 km/s/Mpc)

    # Friedmann equation
    H_squared = H_0**2 * (Omega_m * a**(-3) + Omega_lambda)

    if H_squared < 0:
        return 0

    da_dt = a * np.sqrt(H_squared)
    return da_dt

# Compute black hole properties
masses_bh = np.array([1, 10, 1e6, 4.3e6]) * M_sun  # Solar, stellar, SMBH, Sgr A*
r_schwarzschild = schwarzschild_radius(masses_bh)

# Plot 1: Schwarzschild metric coefficients
fig = plt.figure(figsize=(14, 12))

ax1 = fig.add_subplot(3, 3, 1)
M_test = 10 * M_sun
r_s_test = schwarzschild_radius(M_test)
r_range = np.linspace(1.5 * r_s_test, 10 * r_s_test, 500)

g_tt_vals = -(1 - r_s_test/r_range)
g_rr_vals = 1 / (1 - r_s_test/r_range)

ax1.plot(r_range/r_s_test, -g_tt_vals, 'b-', linewidth=2, label='$-g_{tt}$')
ax1.plot(r_range/r_s_test, g_rr_vals, 'r-', linewidth=2, label='$g_{rr}$')
ax1.axvline(x=1, color='black', linestyle='--', linewidth=1.5, label='Event Horizon')
ax1.set_xlabel('$r/r_s$', fontsize=11)
ax1.set_ylabel('Metric Coefficient', fontsize=11)
ax1.set_title('Schwarzschild Metric Components', fontsize=11)
ax1.legend(fontsize=9)
ax1.grid(True, alpha=0.3)
ax1.set_ylim(0, 5)

# Plot 2: Effective potential comparison
ax2 = fig.add_subplot(3, 3, 2)
M_central = M_sun
L_test = 3e40  # Angular momentum (kg m^2/s)
r_pot = np.linspace(1e9, 5e10, 500)

V_newtonian = effective_potential(r_pot, M_central, L_test, relativistic=False)
V_gr = effective_potential(r_pot, M_central, L_test, relativistic=True)

ax2.plot(r_pot/au, V_newtonian/1e8, 'b--', linewidth=2, label='Newtonian')
ax2.plot(r_pot/au, V_gr/1e8, 'r-', linewidth=2, label='General Relativity')
ax2.set_xlabel('$r$ (AU)', fontsize=11)
ax2.set_ylabel('$V_{eff}/10^8$ (J/kg)', fontsize=11)
ax2.set_title('Effective Potential: GR vs Newtonian', fontsize=11)
ax2.legend(fontsize=9)
ax2.grid(True, alpha=0.3)

# Plot 3: Orbital precession
ax3 = fig.add_subplot(3, 3, 3)

# Solve geodesic equations for precessing orbit
r0 = a_mercury * (1 - e_mercury)  # Perihelion
v0 = np.sqrt(G * M_sun * (1 + e_mercury) / (a_mercury * (1 - e_mercury)))
L_mercury = r0 * v0

y0 = [r0, 0, v0, L_mercury]
t_span = (0, 5 * period_mercury)
t_eval = np.linspace(0, 5 * period_mercury, 5000)

sol = solve_ivp(geodesic_equations, t_span, y0, args=(M_sun,),
                t_eval=t_eval, method='RK45', rtol=1e-8)

r_orbit = sol.y[0]
phi_orbit = sol.y[1]

x_orbit = r_orbit * np.cos(phi_orbit)
y_orbit = r_orbit * np.sin(phi_orbit)

ax3.plot(x_orbit/au, y_orbit/au, 'b-', linewidth=1, alpha=0.7)
ax3.plot(0, 0, 'yo', markersize=12, label='Sun')
ax3.set_xlabel('$x$ (AU)', fontsize=11)
ax3.set_ylabel('$y$ (AU)', fontsize=11)
ax3.set_title('Orbital Precession (Mercury-like)', fontsize=11)
ax3.axis('equal')
ax3.grid(True, alpha=0.3)
ax3.legend(fontsize=9)

# Plot 4: Perihelion precession rates
ax4 = fig.add_subplot(3, 3, 4)

planets = ['Mercury', 'Venus', 'Earth', 'Mars']
semi_major = np.array([57.9, 108.2, 149.6, 227.9]) * 1e9  # meters
eccentricity = np.array([0.2056, 0.0068, 0.0167, 0.0934])

precession_rates = [perihelion_precession_rate(a, e, M_sun)
                   for a, e in zip(semi_major, eccentricity)]

# Convert to arcsec/century
orbits_per_century = 100 * year / (2 * np.pi * np.sqrt(semi_major**3 / (G * M_sun)))
precession_per_century = np.array(precession_rates) * orbits_per_century

ax4.bar(planets, precession_per_century, color='steelblue', edgecolor='black', linewidth=1.5)
ax4.axhline(y=43, color='red', linestyle='--', linewidth=2, label='Mercury Observed')
ax4.set_ylabel('Precession (arcsec/century)', fontsize=11)
ax4.set_title('Perihelion Precession: GR Prediction', fontsize=11)
ax4.legend(fontsize=9)
ax4.grid(True, alpha=0.3, axis='y')

# Plot 5: Light deflection angle
ax5 = fig.add_subplot(3, 3, 5)

# Deflection angle as function of impact parameter
b_range = np.linspace(r_s_test, 100*r_s_test, 500)
deflection_angle = 4 * G * M_test / (c**2 * b_range)  # radians

# Convert to arcseconds
deflection_arcsec = deflection_angle * 180 / np.pi * 3600

ax5.loglog(b_range/r_s_test, deflection_arcsec, 'b-', linewidth=2)
ax5.axhline(y=1.75, color='red', linestyle='--', linewidth=1.5,
           label='Solar limb deflection (1.75\")')
ax5.set_xlabel('Impact Parameter $b/r_s$', fontsize=11)
ax5.set_ylabel('Deflection Angle (arcsec)', fontsize=11)
ax5.set_title('Gravitational Light Deflection', fontsize=11)
ax5.legend(fontsize=9)
ax5.grid(True, alpha=0.3, which='both')

# Plot 6: Gravitational wave chirp
ax6 = fig.add_subplot(3, 3, 6)

M_chirp_gw = 30  # Solar masses (like GW150914)
distance_gw = 410  # Mpc
f_initial_gw = 35  # Hz

t_gw = np.linspace(-0.2, 0, 2000)
h_strain, f_gw = gravitational_wave_strain(t_gw, M_chirp_gw, distance_gw, f_initial_gw)

ax6.plot(t_gw * 1000, h_strain * 1e21, 'b-', linewidth=1)
ax6.set_xlabel('Time to Merger (ms)', fontsize=11)
ax6.set_ylabel('Strain $h \\times 10^{21}$', fontsize=11)
ax6.set_title('Gravitational Wave Chirp Signal', fontsize=11)
ax6.grid(True, alpha=0.3)

# Plot 7: Gravitational wave frequency evolution
ax7 = fig.add_subplot(3, 3, 7)

ax7.plot(t_gw * 1000, f_gw, 'r-', linewidth=2)
ax7.set_xlabel('Time to Merger (ms)', fontsize=11)
ax7.set_ylabel('Frequency (Hz)', fontsize=11)
ax7.set_title('GW Frequency Evolution (Chirp)', fontsize=11)
ax7.grid(True, alpha=0.3)

# Plot 8: FLRW scale factor evolution
ax8 = fig.add_subplot(3, 3, 8)

# Different cosmological models
t_cosmo = np.linspace(0.1, 14, 500)  # Gyr
Omega_m_vals = [1.0, 0.3, 0.0]
Omega_lambda_vals = [0.0, 0.7, 1.0]
labels = ['Matter-only', 'ΛCDM (current)', 'Dark energy only']
colors = ['blue', 'red', 'green']

for Om, Ol, label, color in zip(Omega_m_vals, Omega_lambda_vals, labels, colors):
    a_init = 0.1
    t_sec = t_cosmo * 1e9 * year

    a_solution = odeint(friedmann_equations, a_init, t_sec, args=(Om, Ol))
    ax8.plot(t_cosmo, a_solution.flatten(), linewidth=2, label=label, color=color)

ax8.set_xlabel('Time (Gyr)', fontsize=11)
ax8.set_ylabel('Scale Factor $a(t)$', fontsize=11)
ax8.set_title('Cosmic Expansion: FLRW Models', fontsize=11)
ax8.legend(fontsize=8)
ax8.grid(True, alpha=0.3)

# Plot 9: Event horizon radii
ax9 = fig.add_subplot(3, 3, 9)

mass_labels = ['$1 M_\\odot$', '$10 M_\\odot$', '$10^6 M_\\odot$', 'Sgr A*']
r_s_km = r_schwarzschild / 1000

ax9.barh(mass_labels, r_s_km, color='darkviolet', edgecolor='black', linewidth=1.5)
ax9.set_xlabel('Schwarzschild Radius (km)', fontsize=11)
ax9.set_title('Black Hole Event Horizons', fontsize=11)
ax9.set_xscale('log')
ax9.grid(True, alpha=0.3, which='both', axis='x')

plt.tight_layout()
plt.savefig('general_relativity_analysis.pdf', dpi=150, bbox_inches='tight')
plt.close()
\end{pycode}

\begin{figure}[htbp]
\centering
\includegraphics[width=\textwidth]{general_relativity_analysis.pdf}
\caption{Comprehensive general relativity analysis: (a) Schwarzschild metric coefficients
showing singularity at event horizon $r=r_s$; (b) Effective potential comparing Newtonian
and general relativistic predictions, with GR correction producing stable circular orbits
at different radii; (c) Orbital precession demonstrating non-closed elliptical trajectories;
(d) Perihelion precession rates for inner planets, with Mercury's 43 arcsec/century
matching 1915 observations; (e) Gravitational light deflection as function of impact
parameter; (f-g) Gravitational wave strain and frequency evolution for binary black hole
merger; (h) Cosmic scale factor evolution for different FLRW models; (i) Event horizon
radii for black holes spanning stellar to supermassive scales.}
\label{fig:general_relativity}
\end{figure}

\section{Results}

\subsection{Black Hole Properties}

\begin{pycode}
print(r"\begin{table}[htbp]")
print(r"\centering")
print(r"\caption{Schwarzschild Radii for Representative Black Holes}")
print(r"\begin{tabular}{lccc}")
print(r"\toprule")
print(r"Object & Mass ($M_\odot$) & $r_s$ (km) & $r_s$ (AU) \\")
print(r"\midrule")

for i, label in enumerate(['Solar mass BH', 'Stellar BH', 'IMBH', 'Sgr A*']):
    mass_solar = masses_bh[i] / M_sun
    rs_km = r_schwarzschild[i] / 1000
    rs_au = r_schwarzschild[i] / au

    if mass_solar < 1000:
        print(f"{label} & {mass_solar:.1f} & {rs_km:.2f} & {rs_au:.2e} \\\\")
    else:
        print(f"{label} & {mass_solar:.2e} & {rs_km:.2e} & {rs_au:.3f} \\\\")

print(r"\bottomrule")
print(r"\end{tabular}")
print(r"\label{tab:black_holes}")
print(r"\end{table}")
\end{pycode}

The Schwarzschild radius scales linearly with mass: $r_s = 2GM/c^2 \approx 2.95$ km per
solar mass. For Sagittarius A*, the supermassive black hole at the Milky Way's center with
$M = 4.3 \times 10^6 M_\odot$, the event horizon radius is approximately 0.08 AU, smaller
than Mercury's orbit yet containing millions of solar masses.

\subsection{Orbital Precession Analysis}

\begin{pycode}
print(r"\begin{table}[htbp]")
print(r"\centering")
print(r"\caption{Perihelion Precession Rates for Inner Planets}")
print(r"\begin{tabular}{lcccc}")
print(r"\toprule")
print(r"Planet & $a$ (AU) & $e$ & $\delta\phi$/orbit (\") & $\delta\phi$/century (\") \\")
print(r"\midrule")

for i, planet in enumerate(planets):
    a_au = semi_major[i] / au
    e = eccentricity[i]
    prec_orbit = precession_rates[i]
    prec_cent = precession_per_century[i]
    print(f"{planet} & {a_au:.1f} & {e:.4f} & {prec_orbit:.4f} & {prec_cent:.2f} \\\\")

print(r"\midrule")
print(r"Mercury (observed) & --- & --- & --- & 43.0 \\")
print(r"\bottomrule")
print(r"\end{tabular}")
print(r"\label{tab:precession}")
print(r"\end{table}")
\end{pycode}

Mercury's perihelion precession provided the first precision test of general relativity. The
observed value of 43 arcseconds per century, after accounting for Newtonian perturbations
from other planets, exactly matches Einstein's GR prediction. This $1/r^3$ correction term
in the effective potential arises from the curvature of spacetime near the Sun.

\subsection{Gravitational Wave Properties}

\begin{pycode}
# Calculate properties for GW150914-like event
M1_gw150914 = 36  # Solar masses
M2_gw150914 = 29
M_chirp_150914 = (M1_gw150914 * M2_gw150914)**(3/5) / (M1_gw150914 + M2_gw150914)**(1/5)
distance_150914 = 410  # Mpc
peak_strain = 1.0e-21

# Frequency at merger
f_merge = c**3 / (6**(3/2) * np.pi * G * (M1_gw150914 + M2_gw150914) * M_sun)

print(r"\begin{table}[htbp]")
print(r"\centering")
print(r"\caption{Gravitational Wave Event Parameters (GW150914-type)}")
print(r"\begin{tabular}{lc}")
print(r"\toprule")
print(r"Parameter & Value \\")
print(r"\midrule")
print(f"Chirp mass $\\mathcal{{M}}$ & {M_chirp_150914:.1f} $M_\\odot$ \\\\")
print(f"Luminosity distance & {distance_150914} Mpc \\\\")
print(f"Peak strain amplitude & {peak_strain:.1e} \\\\")
print(f"Merger frequency & {f_merge:.0f} Hz \\\\")
print(f"Total radiated energy & $\\sim 3 M_\\odot c^2$ \\\\")
print(r"\bottomrule")
print(r"\end{tabular}")
print(r"\label{tab:gw_event}")
print(r"\end{table}")
\end{pycode}

\section{Discussion}

\begin{example}[Testing General Relativity]
GR has passed precision tests across vast scales:
\begin{itemize}
\item \textbf{Solar System}: Mercury precession (43\"/century), light deflection by Sun (1.75\")
\item \textbf{Binary Pulsars}: Orbital decay matching gravitational wave emission (Hulse-Taylor)
\item \textbf{Cosmology}: Accelerating expansion requiring dark energy ($\Omega_\Lambda \approx 0.7$)
\item \textbf{Strong Gravity}: LIGO/Virgo detections confirming black hole mergers
\end{itemize}
\end{example}

\begin{remark}[Gravitational Time Dilation]
The Schwarzschild metric predicts time runs slower in gravitational fields. For GPS satellites
at altitude $h \approx 20{,}000$ km, clocks run approximately 38 microseconds per day faster
than Earth-surface clocks due to weaker gravity (offset by 7 $\mu$s/day special relativistic
correction). Without GR corrections, GPS would accumulate 10 km positioning errors daily.
\end{remark}

\subsection{Event Horizon Physics}

\begin{pycode}
# Calculate tidal forces at horizon
M_bh_tidal = 10 * M_sun  # Stellar mass black hole
r_s_tidal = schwarzschild_radius(M_bh_tidal)

# Tidal acceleration for 2m tall human
height = 2.0  # meters
tidal_accel = 2 * G * M_bh_tidal * height / r_s_tidal**3

print(f"For a {M_bh_tidal/M_sun:.0f} $M_\\odot$ black hole with $r_s = {r_s_tidal/1000:.1f}$ km, ")
print(f"tidal acceleration across a {height} m object at the event horizon is ")
print(f"${tidal_accel:.2e}$ m/s$^2$ ($\\sim {tidal_accel/9.8:.1e}$ $g$), ")
print(f"demonstrating extreme spaghettification effects.")
\end{pycode}

\subsection{Cosmological Implications}

The FLRW metric with $\Omega_m = 0.3$ and $\Omega_\Lambda = 0.7$ (current cosmological
observations) predicts eternal accelerating expansion. The scale factor grows exponentially
at late times: $a(t) \propto e^{H_0\sqrt{\Omega_\Lambda}\,t}$, with galaxies beyond the
cosmic event horizon ($\sim 16$ Gpc) forever receding faster than light, rendering them
causally disconnected from future observers.

\section{Conclusions}

This computational analysis demonstrates general relativity's fundamental predictions:

\begin{enumerate}
\item Schwarzschild black holes exhibit event horizons at $r_s = 2GM/c^2$, ranging from 3 km for
solar-mass objects to $10^{10}$ km for supermassive black holes
\item Orbital precession arises from the $-GML^2/(c^2r^3)$ potential correction, yielding
Mercury's observed 43 arcseconds per century perihelion advance
\item Gravitational waves propagate as ripples in spacetime, with LIGO-detected strains
$h \sim 10^{-21}$ confirming binary black hole mergers
\item FLRW cosmology with $\Lambda > 0$ describes our accelerating universe, with current
observations favoring $\Omega_\Lambda \approx 0.7$
\item Light deflection, time dilation, and geodesic motion provide precision tests confirming
GR from Solar System scales to cosmological horizons
\end{enumerate}

Einstein's geometric theory of gravity remains our most accurate description of spacetime,
validated across 60 orders of magnitude in length scale.

\section*{Further Reading}

\begin{itemize}
\item Misner, C.W., Thorne, K.S., \& Wheeler, J.A. \textit{Gravitation}. Freeman, 1973.
\item Weinberg, S. \textit{Gravitation and Cosmology}. Wiley, 1972.
\item Carroll, S.M. \textit{Spacetime and Geometry: An Introduction to General Relativity}. Cambridge, 2019.
\item Wald, R.M. \textit{General Relativity}. University of Chicago Press, 1984.
\item Hartle, J.B. \textit{Gravity: An Introduction to Einstein's General Relativity}. Addison-Wesley, 2003.
\item Schutz, B.F. \textit{A First Course in General Relativity}, 3rd ed. Cambridge, 2022.
\item Abbott, B.P., et al. (LIGO/Virgo). ``Observation of Gravitational Waves from a Binary Black Hole Merger.'' \textit{Phys. Rev. Lett.} \textbf{116}, 061102 (2016).
\item Will, C.M. ``The Confrontation between General Relativity and Experiment.'' \textit{Living Rev. Relativity} \textbf{17}, 4 (2014).
\item Perlmutter, S., et al. ``Measurements of $\Omega$ and $\Lambda$ from 42 High-Redshift Supernovae.'' \textit{Astrophys. J.} \textbf{517}, 565 (1999).
\item Event Horizon Telescope Collaboration. ``First M87 Event Horizon Telescope Results. I. The Shadow of the Supermassive Black Hole.'' \textit{Astrophys. J. Lett.} \textbf{875}, L1 (2019).
\item Riess, A.G., et al. ``Observational Evidence from Supernovae for an Accelerating Universe.'' \textit{Astron. J.} \textbf{116}, 1009 (1998).
\item Stairs, I.H. ``Testing General Relativity with Pulsar Timing.'' \textit{Living Rev. Relativity} \textbf{6}, 5 (2003).
\item Hulse, R.A., \& Taylor, J.H. ``Discovery of a Pulsar in a Binary System.'' \textit{Astrophys. J. Lett.} \textbf{195}, L51 (1975).
\item Ashby, N. ``Relativity in the Global Positioning System.'' \textit{Living Rev. Relativity} \textbf{6}, 1 (2003).
\item Shapiro, I.I. ``Fourth Test of General Relativity.'' \textit{Phys. Rev. Lett.} \textbf{13}, 789 (1964).
\end{itemize}

\end{document}
